% ===== Document Preamble =====
\documentclass[11pt]{article}
\usepackage[margin=0.5in]{geometry}
\usepackage{amsmath,amssymb}
\usepackage{xcolor}
\usepackage{enumitem}
\usepackage{needspace}
\usepackage{comment}

% ===== Colors =====
\definecolor{AnsRed}{RGB}{170,0,0}
\definecolor{WrongBlue}{RGB}{0,0,160}
\definecolor{SkipGray}{RGB}{120,120,120}

% ===== Choice list styling (tight, consistent) =====
\setlist[enumerate]{%
  leftmargin=1.5em,
  itemsep=0.25em,
  topsep=0.25em,
  label=\textbf{\Alph*.}
}

% ===== Style Wrappers =====
\newcommand{\QuestionStyle}[1]{{\large\bfseries\boldmath #1}}
\newcommand{\AnswerStyle}[1]{{\large\bfseries\boldmath #1}}
\newcommand{\ReasonStyle}[1]{{\normalsize #1}}

% ===== Show / Hide Reasoning Toggle =====
\newif\ifshowreason
\showreasontrue      % Instructor version
% \showreasonfalse   % Student version

% ===== Choice Item Macros (Answer + Reasoning) =====
% Usage:
%   \Correct{Answer}{Reasoning}
%   \Wrong{Answer}{Reasoning}

\newcommand{\Correct}[2]{%
  \item \AnswerStyle{\textcolor{AnsRed}{#1}}%
  \ifshowreason\\[-0.2em]
    \ReasonStyle{\textcolor{AnsRed}{#2}}%
  \fi
}

\newcommand{\Wrong}[2]{%
  \item \AnswerStyle{\textcolor{WrongBlue}{#1}}%
  \ifshowreason\\[-0.2em]
    \ReasonStyle{\textcolor{WrongBlue}{#2}}%
  \fi
}

% ===== Multiple-Choice Question Block Macro =====
% Usage: \QBlock{Q1}{stem...}{<items>}{answer key}

\newcommand{\QBlock}[4]{%
  \Needspace{8\baselineskip}%
  \noindent
  \begin{minipage}{\linewidth}
    \QuestionStyle{#1.~#2}

    \vspace{0.25em}
    \begin{enumerate}
      #3
    \end{enumerate}

    \vspace{0.1em}
    \noindent\textcolor{AnsRed}{\textit{\textbf{Answer:}~#4}}

    \vspace{0.25em}
    \hrule height 0.4pt
  \end{minipage}
  \vspace{0.75em}
}

% ===== Open-Ended Question Block Macro =====
% Usage: \OEBlock{Q1}{stem...}{answer text}

\newcommand{\OEBlock}[3]{%
  \Needspace{5\baselineskip}%
  \noindent
  \begin{minipage}{\linewidth}
    \QuestionStyle{#1.~#2}

    \vspace{0.25em}
    \noindent\textcolor{AnsRed}{\textit{\textbf{Answer:}~#3}}

    \vspace{0.25em}
    \hrule height 0.4pt
  \end{minipage}
  \vspace{0.75em}
}

% ===== Multi-Part Question Block Macro =====
% Usage: \MPBlock{Q1}{stem...}{<parts>}

\newcommand{\MPBlock}[3]{%
  \Needspace{8\baselineskip}%
  \noindent
  \begin{minipage}{\linewidth}
    \QuestionStyle{#1.~#2}

    \vspace{0.25em}
    #3

    \vspace{0.25em}
    \hrule height 0.4pt
  \end{minipage}
  \vspace{0.75em}
}

% ===== MC Sub-Part Macro =====
% Usage: \MCPart{a}{stem...}{<items>}{answer key}

\newcommand{\MCPart}[4]{%
  \par\vspace{0.35em}
  \noindent\textbf{(#1)}~#2
  \begin{enumerate}
    #3
  \end{enumerate}
  \vspace{0.1em}
  \noindent\textit{\textbf{Answer:}~#4}
  \par
}

% ===== Open-Ended Sub-Part Macro =====
% Usage: \OEPart{a}{stem...}{answer text}

\newcommand{\OEPart}[3]{%
  \par\vspace{0.35em}
  \noindent\textbf{(#1)}~#2
  \par\vspace{0.15em}
  \noindent\textit{\textbf{Answer:}~#3}
  \par
}

% ===== Skipped Question Note Macro =====
% Usage: \SkipBlock{Q1}{reason}

\newcommand{\SkipBlock}[2]{%
  \noindent\textcolor{SkipGray}{\textit{[#1 omitted: #2]}}%
  \vspace{0.5em}\par
}

% ===== Document Start =====

\begin{document}

\begin{center}
{\huge \textbf{Quantum Mechanics: Wave Functions and the Infinite Square Well}}\\[0.25em]
{\Large Isomorphic Practice Questions}\\[0.25em]
{\large Questions generated below. Multiple-choice answers in RED; incorrect answers in BLUE.}
\end{center}

\vspace{0.75em}

\MPBlock{Q1}{A particle of mass $m$ in one dimension has the wave function
\[
\Psi(x,t)=C\,\exp\!\left[-\lambda\left(\frac{m x^2}{2\hbar}+ i t\right)\right],\qquad \lambda>0.
\]
}{
  \OEPart{a}{Find the normalization constant $C$ (take it to be real and positive).}{Normalize: $1=\int_{-\infty}^{\infty}|\Psi(x,t)|^2\,dx=|C|^2\int_{-\infty}^{\infty}\exp\!\left(-\lambda\frac{m x^2}{\hbar}\right)dx=|C|^2\sqrt{\frac{\pi\hbar}{\lambda m}}$. Hence $|C|^2=\sqrt{\frac{\lambda m}{\pi\hbar}}$ and one convenient choice is $C=\left(\frac{\lambda m}{\pi\hbar}\right)^{1/4}$.}
  \OEPart{b}{Write the integral you would evaluate to compute $\langle x\rangle$ at time $t$, and determine $\langle x\rangle$.}{Use $\langle x\rangle=\int_{-\infty}^{\infty}\Psi^*(x,t)\,x\,\Psi(x,t)\,dx=\int_{-\infty}^{\infty}x\,|\Psi(x,t)|^2\,dx$. Here $|\Psi|^2=|C|^2\exp\!\left(-\lambda\frac{m x^2}{\hbar}\right)$ is even in $x$, so $x|\Psi|^2$ is odd and integrates to $0$; therefore $\langle x\rangle=0$.}
}

\OEBlock{Q2}{At time $t=0$, a particle is described by the (not necessarily normalized) wave function
\[
\Psi(x,0)=
\begin{cases}
K\,\dfrac{x}{p}, & 0\le x\le p,\\[6pt]
K\,\dfrac{q-x}{q-p}, & p\le x\le q,\\[6pt]
0, & \text{otherwise},
\end{cases}
\]
where $p$ and $q$ are positive constants with $0<p<q$, and $K=\dfrac{5}{2q}$. What is the probability of finding the particle to the left of $x=p$?}{The probability is $P(x<p)=\int_{-\infty}^{p}|\Psi(x,0)|^2dx=\int_{0}^{p}\left|K\frac{x}{p}\right|^2dx=\frac{K^2}{p^2}\int_{0}^{p}x^2dx=\frac{K^2 p}{3}$. With $K=\frac{5}{2q}$, this gives $P(x<p)=\frac{25}{4q^2}\cdot\frac{p}{3}=\frac{25p}{12q^2}$.}

\OEBlock{Q3}{A state is given by
\[
\Phi(x,t)=\varphi(x)\,e^{-iEt/\hbar},
\]
where $E$ is a constant. Using the probability density, explain why this represents a stationary state.}{A stationary state has a time-independent probability density. Here
\[
|\Phi(x,t)|^2=\Phi^*(x,t)\Phi(x,t)=|\varphi(x)|^2\,\big|e^{-iEt/\hbar}\big|^2=|\varphi(x)|^2,
\]
so the time dependence is only a global phase and does not change any probabilities.}

\MPBlock{Q4}{A normalized wavefunction is expanded in energy eigenstates as
\[
\Psi(x,t)=\sum_{n=1}^{\infty} d_n\,u_n(x)\,e^{-i\varepsilon_n t/\hbar},
\]
where $H u_n=\varepsilon_n u_n$ and the $u_n$ form an orthonormal set.}{
  \OEPart{a}{What is the physical meaning of $|d_n|^2$?}{If the state is normalized and expanded in energy eigenstates, then $|d_n|^2$ is the probability that an energy measurement yields the value $\varepsilon_n$. Equivalently, it is the population (weight) of the $n$th stationary state in the superposition.}
  \OEPart{b}{Write the expectation value of the energy, $\langle H\rangle$, in terms of $|d_n|^2$ and $\varepsilon_n$.}{For an expansion in energy eigenstates, the energy expectation value is
\[
\langle H\rangle=\sum_{n=1}^{\infty} |d_n|^2\,\varepsilon_n,
\]
using $\sum_n |d_n|^2=1$ for a normalized state. (If not normalized, then $\langle H\rangle=\dfrac{\sum_n |d_n|^2\varepsilon_n}{\sum_n |d_n|^2}$.)}
}

\MPBlock{Q5}{Consider a particle of mass $m$ in a one-dimensional infinite square well of width $L$:
\[
V(x)=\begin{cases}
0,& 0\le x\le L,\\
\infty,& \text{otherwise}.
\end{cases}
\]
}{
  \OEPart{a}{Using $k^2=2mE/\hbar^2$, write the time-independent Schr\"odinger equation inside the well and give its general solution for $0<x<L$.}{Inside the well ($V=0$), the TISE is $-\frac{\hbar^2}{2m}\,\frac{d^2\psi}{dx^2}=E\psi$, i.e.
\[\frac{d^2\psi}{dx^2}+k^2\psi=0,\qquad k^2=\frac{2mE}{\hbar^2}.
\]
The general solution for $0<x<L$ is $\psi(x)=A\sin(kx)+B\cos(kx)$, and $\psi=0$ outside since $V=\infty$.}
  \OEPart{b}{Apply the boundary conditions to determine the allowed energies $E_n$.}{Infinite walls impose $\psi(0)=\psi(L)=0$. From $\psi(0)=0$ we get $B=0$, and $\psi(L)=0$ gives $\sin(kL)=0\Rightarrow kL=n\pi$ with $n=1,2,\dots$. Thus $k_n=\frac{n\pi}{L}$ and
\[
E_n=\frac{\hbar^2 k_n^2}{2m}=\frac{n^2\pi^2\hbar^2}{2mL^2}.
\]}
  \OEPart{c}{Find the normalized energy eigenfunctions $\psi_n(x)$.}{With $B=0$ and $k_n=n\pi/L$, the eigenfunctions are $\psi_n(x)=A_n\sin\left(\frac{n\pi x}{L}\right)$ for $0<x<L$ and $0$ otherwise. Normalize:
\[
1=\int_0^L |\psi_n(x)|^2dx=|A_n|^2\int_0^L\sin^2\left(\frac{n\pi x}{L}\right)dx=|A_n|^2\frac{L}{2},
\]
so $A_n=\sqrt{\frac{2}{L}}$. Hence $\psi_n(x)=\sqrt{\frac{2}{L}}\sin\left(\frac{n\pi x}{L}\right)$ in the well.}
  \OEPart{d}{Show that $\psi_n$ and $\psi_m$ are orthogonal for $n\neq m$.}{Using the normalized functions,
\[
\int_0^L \psi_n(x)\psi_m(x)\,dx=\frac{2}{L}\int_0^L \sin\left(\frac{n\pi x}{L}\right)\sin\left(\frac{m\pi x}{L}\right)dx.
\]
The sine orthogonality integral is $0$ for $n\neq m$, so $\int_0^L \psi_n\psi_m\,dx=0$, proving orthogonality.}
}

\OEBlock{Q6}{In a 1D infinite square well, suppose the particle's initial state is a superposition of the first two stationary states:
\[
\Psi(x,0)=\sqrt{\frac{2}{5}}\,\psi_1(x)+\sqrt{\frac{3}{5}}\,\psi_2(x).
\]
Find the expectation value of the energy $\langle E\rangle$ in terms of $E_1$ and $E_2$ (you may assume the $\psi_n$ are orthonormal energy eigenstates).}{For a superposition of energy eigenstates, $\langle E\rangle=\sum_n |c_n|^2E_n$. Here $c_1=\sqrt{\tfrac{2}{5}}$ and $c_2=\sqrt{\tfrac{3}{5}}$, so
\[
\langle E\rangle=\frac{2}{5}E_1+\frac{3}{5}E_2.
\]}

\end{document}
